\documentclass{article}
\usepackage{style}

\title{Experiment 7 - Bragg's Law}
\author{Daire O'Brien - 24365676}
\begin{document}
    \maketitle
    \section{Background}
        In 1912 German physicist Max von Laue discovered that when X-rays passed through a crystal sample diffraction and interferance occured. The effect was observed when a crystal was struck with a broad-band X-ray source and manifested as a ring-like pattern of dots on a photographic plate placed behind the crystal. \img{LauDiffractionDiagram}{Interferance pattern from X-rays passing through a crystal sample}{X}{0.35} As von Laue was investigating if X-rays were particles or waves, physicists William H. Bragg and William L. Bragg investigated if the crystals could be used to reflect the X-rays. \par
        Using an x-ray source, an adjustable apature and a turntable they were able to determine that scattering only occured when there was constructive interferance between the x-rays. They proposed that 
            \begin{equation}
                n \lambda = 2d \sin \theta
            \end{equation}
        was the condition for the  reflected x-rays to be detected. \par
        More recently, 
    \section{References}
        \begin{enumerate}
            \item William Henry Bragg, William Lawrence Bragg, 1913, \emph{The reflection of X-rays by crystals}, Proc. R. Soc. A1 July 1913
            \item  W. Friedrich, P. Knipping, M. Laue, presented by A. Sommerfeld 1912, \emph{Interference Phenomena for X-Rays, Theoretical section}, Proceedings of the Mathematical-Physical Class of the Royal Bavarian Academy of Sciences in Munich, Issue 2. English translation by J. J. Stezowski.
        \end{enumerate}
\end{document}